% One addition saved by forcing the first Pell index to be a multiple.
\section{Forcing a multiple of the Pell index}
\label{sec:index-multiple}

The first Pell index can be isolated with a single product in place of
an affine congruence. Arrange that its parameter is congruent to one
modulo the intended index. Then require its denominator to be a
positive multiple of that index. Pell growth excludes every multiple
after the first.

Retain the fixed encoding and the shifted main parameter from
Section~\ref{sec:main-base-four}, but change the intermediate scale to
\[
\begin{gathered}
 N=n,\quad R=r,\quad U=wN^2,\quad Y_P=sN^2,\quad D=R+1,\\
 M=DY_P,\quad a_0=M(U+1),\quad A=a_0+4,
 \quad Q=UM^2,\quad P=2Q+1,\quad J=2R+1.
\end{gathered}
\]
The supplied unknown $a$ still denotes $a_0$. Replace the old equation
$k=R+1+hQ$ by
\begin{equation}\label{eq:index-multiple}
 k=h(R+1)=hD,\qquad h>0.
\end{equation}
Retain the triangular norm and positive interval
\begin{equation}\label{eq:index-multiple-norm}
 \tau(\tau+1)=Q(Q+1)k^2,\qquad
 c=kY_P+\eta,\qquad k=\eta+\zeta,
 \qquad \tau,\eta,\zeta>0.
\end{equation}
The scale $M$ is an intermediate value, so the changed scale affects
the expanded source polynomials E9 and E12; equation
\eqref{eq:index-multiple} changes E11. All other source polynomials
in the supplied variables are unchanged.

\subsection{Isolating the first positive multiple}

The unchanged pre-Pell coding proof supplies
\[
 N=q^8,\quad q>4b\ge8,\quad
 3L\le B\le N\le R,\quad b,q\le N.
\]
Thus $N,R\ge64$, $U,Y_P\ge4096$, and $M\ge262144$.
Equation~\eqref{eq:index-multiple} gives $k\ge R+1$, and the interval
in~\eqref{eq:index-multiple-norm} gives
\[
 c>Y_Pk\ge Y_P(R+1)>2R+1=J.
\]
The generic signed Pell argument therefore gives
$c=\psi_A(J)$ and $d=\chi_A(J)$.

Put $\widehat\tau=2\tau+1$. The norm in
\eqref{eq:index-multiple-norm} is precisely
\[
 \widehat\tau^2-(P^2-1)k^2=1,\qquad \widehat\tau>1.
\]
Hence $k=\psi_P(t)$ for a positive integer $t$.
The new scale makes
\[
 P-1=2Q=2U(R+1)^2Y_P^2
\]
divisible by $D$. The elementary Pell congruence
$\psi_P(t)\equiv t\pmod{P-1}$ therefore gives
$k\equiv t\pmod D$. Equation~\eqref{eq:index-multiple} yields
\begin{equation}\label{eq:index-positive-multiple}
 t=mD=m(R+1),\qquad m\ge1.
\end{equation}

If $m\ge2$, then $t-1\ge2R+1$, and Pell growth gives
\[
 \frac ck\le\frac{(2A)^{2R}}{(2P-1)^{2R+1}}
 <\left(\frac{2A}{2P-1}\right)^{2R}<\frac12.
\]
The last inequality follows from $2A/(2P-1)<1/2$; with the shifted
parameter this is equivalent to
\[
 4UM(M-1)-4M-15=4M\bigl(U(M-1)-1\bigr)-15>0.
\]
The initial bounds already imply it. This contradicts $c/k>Y_P$.
Consequently
\begin{equation}\label{eq:index-multiple-exact}
 k=\psi_P(R+1),\qquad\widehat\tau=\chi_P(R+1).
\end{equation}
The conclusion uses no exponent relation or binomial rounding.
No parity condition on the positive quotient $h$ is required.

\subsection{Preserving the approximation and the encoding}

At the recovered indices, the principal ratio still cancels the scale:
\[
 \frac{(2M(U+1))^{2R}}{(4UM^2)^R}
 =\xi:=\frac{(U+1)^{2R}}{U^R}.
\]
The lower estimate of Section~\ref{sec:main-base-four} uses only
$14Q>a_0$, which remains true. Its upper estimate uses
\[
 \frac{8R}{a_0}
 =\frac{8R}{(R+1)Y_P(U+1)}
 <\frac8{Y_P(U+1)}<\frac12.
\]
Thus the same growth and binomial estimates give
\begin{equation}\label{eq:index-multiple-ratio}
 \xi<\frac ck
 <\xi\left(1+\frac{16R}{a_0}\right)
 <\xi\left(1+\frac{16}{Y_P(U+1)}\right).
\end{equation}
In particular $\xi<Y_P+1$, so $Y_P\ge U^R$ and
\[
 A>a_0=(R+1)Y_P(U+1)>RU^R,
 \qquad 0<\frac ck-\xi<\frac{32}{U+1}<\frac12.
\]
These are exactly the size and error bounds needed after index
recovery in Section~\ref{sec:main-base-four}; they have been
established here for the changed scale.

Explicitly, $N,R\ge64$ and $bw\le U$ give
\[
 4^{3J}\le64N^{2R}\le RU^R<A,
 \qquad (bw)^3\le U^R<A.
\]
The unchanged base-four congruence yields $bw=4^J$, so $b$ is a
power of two and $U\ge N4^J>4\cdot2^{2R}$.
The unchanged second Pell norm, gap and index congruence give
$\kappa=\psi_A(L)$ and $\mu=\chi_A(L)$.
The inequalities
\[
 B^{3L}\le B^B\le N^N\le U^R<A,
 \qquad q^3\le N^N\le U^R<A
\]
then give $q=B^L$ through the second exponent congruence.
For $F=\lfloor\xi\rfloor$, the binomial expansion gives
\[
 0<\xi-F<\frac14,\qquad
 F\equiv\binom{2R}{R}\pmod U.
\]
Now $F<\xi<c/k<Y_P+1$ and
$Y_P<c/k<\xi+1/2<F+3/4$ imply $Y_P=F$ by integrality.
Thus $N^2\mid\binom{2R}{R}$. The unchanged no-carry and
modular-Sidon decoding proves membership in the represented system.

\subsection{Positive witnesses and the saved addition}

Conversely, start with the canonical coding witnesses for a solution
of the represented system. They determine $b,B,q,N,R$ before any Pell
coordinates are chosen. As in Section~\ref{sec:main-base-four}, put
\[
 J=2R+1,\qquad w=4^J/b,\qquad U=N^2w,
 \qquad Y_P=\lfloor\xi\rfloor,\qquad s=Y_P/N^2.
\]
The positive integrality of $w$ and $s$ depends only on these coding
values and the binomial expansion, so it is unaffected by the scale
change. Now form $M=(R+1)Y_P$, $a_0=M(U+1)$, $A=a_0+4$,
$Q=UM^2$, $P=2Q+1$, and choose
\[
 c=\psi_A(J),\qquad d=\chi_A(J),\qquad k=\psi_P(R+1).
\]
Estimate~\eqref{eq:index-multiple-ratio} at these chosen indices
gives $Y_P<c/k<Y_P+3/4$, so
$\eta=c-kY_P$ and $\zeta=k-\eta$ are positive integers.
Choose
\[
 \tau=\frac{\chi_P(R+1)-1}{2},\qquad h=\frac{k}{R+1}.
\]
The oddness of $P$ makes $\tau$ positive integral. Since
$R+1\mid P-1$ and $\psi_P(R+1)\equiv R+1\pmod{P-1}$,
$h$ is a positive integer and satisfies the new product equation.
It is a supplied positive unknown; no division operation is omitted
from the certificate.

Choose the remaining dependent Pell witnesses using the generic
constructions at this new $A$ from Section~\ref{sec:main-base-four}.
In particular, after choosing positive $f,i$ for E16, take
\[
 G=1+(A+1)(f^2-1),\qquad
 o=\frac{\chi_G(J)+d}{f},\qquad
 j=\frac{\psi_G(J)-J}{c},
\]
and take $\kappa=\psi_A(L)$, $\mu=\chi_A(L)$,
$\Delta=(\kappa-L)/(A-1)$ and $\phi=c-\kappa$.
The same generic congruences and strict Pell growth make them positive
integers. The exponent quotients are positive because
\[
 d-(A-4)c>3c>4^J,\qquad
 \mu-(A-B)\kappa>(B-1)\kappa>B^L.
\]
Their moduli are positive, and the required inequalities
$A>(4^J)^3,(B^L)^3$ were preserved above. Thus every witness is
constructed with its required positive domain. All coding witnesses
remain valid, while $M$ and its dependent Pell coordinates are
recomputed; the old and new values of $h$ are not identified.

\begin{theorem}\label{thm:best105}
For each r.e.\ set, the fixed modular-Sidon encoding gives a
membership-equivalent system in 34 positive unknowns and 22 equations
with a straight-line certificate of \textbf{105 arithmetic instructions}:
59 multiplications and 46 additions. Fixed numerals, supplied parameters
and equality tests are free.
\end{theorem}
\begin{proof}
The preceding arguments prove soundness and construct all positive
witnesses for necessity. For the count, the existing instruction
$r_1=r+1$ is moved before the computation of $M$. Replacing
$M=rY_P$ by $M=r_1Y_P$ retains one multiplication. Replacing $hQ$
by $hr_1$ also retains one multiplication, but the old addition
$r_1+hQ$ disappears: the free equality now compares $k$ directly
with $hr_1$.

The other instructions retain their counts. Hence the 106-operation
certificate loses one addition and has 59 multiplications and 46
additions. The positive unknown and equation counts stay unchanged.
The checker \file{round16\_1980\_index\_multiple\_certificate.py}
verifies the three changed source residuals, every unchanged residual,
both triangular corrections, all modified register identities, and
the full serialized primitive certificate. Its fixed exponent is
$L=2^{32}$.
\end{proof}
